(1) 旋转盘上半径为 $r$ 的导体质元速度为 $\boldsymbol v=\Omega r\hat{\boldsymbol\phi}$，洛伦兹力的正电荷方向为 $\boldsymbol v\times\boldsymbol B=\Omega rB\hat{\boldsymbol r}$。沿半径从圆心积分动生电动势：
\[
\mathcal E=\int_0^a(\boldsymbol v\times\boldsymbol B)\cdot d\boldsymbol\ell
=\int_0^a B\Omega r\,dr=\frac12B\Omega a^2.
\]
盘外缘电势高于圆心，闭合回路中正电荷电流沿盘由内向外，外负载再回到圆心。忽略电感时 $I=\mathcal E/R=B\Omega a^2/(2R)$。

(2) 盘内径向电流元 $I\,dr\,\hat r$ 与 $\boldsymbol B$ 的叉乘沿 $-\hat{\boldsymbol\phi}$，每元所受电磁阻力矩大小为 $d\tau=IBr\,dr$。积分得
\[
\tau_B=-\frac12IBa^2=-\frac{B^2a^4}{4R}\Omega .
\]
令 $\gamma_B=B^2a^4/(4R)>0$，则转动方程
\[
J\dot\Omega=\tau_0-\gamma_B\Omega .
\]
电磁转矩阻碍旋转，符合楞次定律；若错误地取正号会得到无需外界能量的自发加速。

(3) 初角速度零时上述线性方程的精确解是
\[
\Omega(t)=\Omega_\infty(1-e^{-t/\tau}),\qquad
\Omega_\infty=\frac{\tau_0}{\gamma_B}=\frac{4R\tau_0}{B^2a^4},
\qquad\tau=\frac J{\gamma_B}.
\]
电流 $I(t)=Ba^2\Omega(t)/(2R)$。终端状态电磁阻力矩正好抵消电机转矩，$\dot\Omega\to0$，但盘与外电阻仍持续有能量流入及耗散。

(4) 对方程乘 $\Omega$：
\[
\tau_0\Omega=\frac{d}{dt}\left(\frac J2\Omega^2\right)+\gamma_B\Omega^2.
\]
又电阻热功率 $I^2R=(B^2a^4/(4R))\Omega^2=\gamma_B\Omega^2$，两种推导完全吻合。若固定有限时间令 $R\to\infty$，$\gamma_B\to0$、$\tau\to\infty$，圆盘近似按 $J\dot\Omega=\tau_0$ 匀加速，闭合电流和磁制动力矩趋零；先取每个有限 $R$ 的 $t\to\infty$ 则得随 $R$ 增大的终端速度，两种极限不能交换。
